% ---------------- Math macros ----------------
\newcommand{\R}{\mathbb{R}}
\newcommand{\N}{\mathbb{N}}
\newcommand{\Z}{\mathbb{Z}}
\newcommand{\E}{\mathbb{E}}
\renewcommand{\Pr}{\mathbb{P}}
% \O already exists in LaTeX (the letter O-slash), so keep it for text mode
% and use \mathcal{O} (Big-O) in math mode.
\let\oslashletter\O
\renewcommand{\O}{\ifmmode\mathcal{O}\else\oslashletter\fi}
\newcommand{\norm}[1]{\left\lVert #1 \right\rVert}
\newcommand{\abs}[1]{\left\lvert #1 \right\rvert}
\newcommand{\inner}[2]{\left\langle #1, #2 \right\rangle}
\DeclareMathOperator*{\argmin}{arg\,min}
\DeclareMathOperator*{\argmax}{arg\,max}

% Nicer QED symbol
\renewcommand{\qedsymbol}{$\blacksquare$}

% Inline code:  \code{merge\_sort(A)}
\newcommand{\code}[1]{\colorbox{codebg}{\texttt{\small #1}}}

\newcommand{\grupointegrantes}[2]{%
  \begin{tcolorbox}[colback=white!92!accent, colbacktitle=accent,
    colframe=accent, coltitle=white, coltext=black,
    title={\sffamily\bfseries #1 --- Integrantes},
    boxrule=0.6pt, arc=1pt, left=14pt, right=14pt, top=6pt, bottom=6pt,
    toptitle=4pt, bottomtitle=4pt]
    \small
    \begin{tabularx}{\linewidth}{@{}>{\raggedright\arraybackslash}X r
                                    >{\raggedright\arraybackslash}X r@{}}
      #2
    \end{tabularx}
  \end{tcolorbox}%
  \vspace{8pt}%
}
